\chapter{Imaginary Time, Poisson, and the Free Boundary}\label{ch:relaxation}
\index{relaxation dynamics}\index{imaginary-time propagation}\index{Poisson equation}\index{level-set method}\index{three-line code}

The Euler--Lagrange equations of Chapter~\ref{ch:equation} are a coupled system: each $\psi_i$ satisfies a Schr\"odinger eigenvalue problem on its own domain $\Omega_i$, the domains themselves satisfy a Bernoulli free-boundary condition, and the Hartree potentials $V_j$ couple every electron to every other. There is no closed-form solution, but there is a remarkably clean way to relax the system to its minimum --- a method that consists, at its mathematical core, of just three coupled time-evolution equations applied at every grid point and every time step.

We refer to this as the \textbf{three-line code}: three lines of PDE that, applied iteratively, drive any initial configuration toward the variational minimum of the RealQM energy. Everything else in the implementation --- the discretisation, the GPU dispatch, the diagnostics, the molecular dynamics --- is data plumbing around those three lines. This chapter writes the three lines down, explains what each one is doing physically, and describes how the partition is updated together with the wave functions.

\section{The three-line code}\label{sec:three-line}

For $i = 1, \ldots, N$, the relaxation evolves three quantities simultaneously: the wave function $\psi_i$, the per-electron Hartree potential $V_i$, and an indicator $\chi_i$ of the domain $\Omega_i$. The three coupled equations are
\begin{align}
\dot{\psi}_i + \mathcal{H}_i\, \psi_i &= 0 \quad\text{in } \Omega_i, \quad \frac{\partial \psi_i}{\partial n} = 0 \quad\text{on } \Gamma_i, \label{eq:3line-psi} \\
\dot{V}_i - \Delta V_i &= 2\pi\, \psi_i^2 \quad\text{in } \mathbb{R}^3, \label{eq:3line-V} \\
\dot{\chi}_i + \beta_i\, |\nabla \chi_i| &= 0 \quad\text{in } \mathbb{R}^3, \label{eq:3line-chi}
\end{align}
with $\mathcal{H}_i = -\tfrac{1}{2}\Delta + W + 2\sum_{j \ne i} V_j$ from~\eqref{eq:Hi-def}, and $\psi_i$ renormalised to $\int \psi_i^2 = 1$ at every step. Each equation corresponds to one of the three physical processes the framework needs.

\paragraph{\eqref{eq:3line-psi} is imaginary-time propagation.} The eigenvalue problem~\eqref{eq:Hi} for the ground state of $\mathcal{H}_i$ can be solved by evolving $\psi_i$ in \emph{imaginary time} $\tau = -i t$:
\[
\dot{\psi}_i \;=\; -\mathcal{H}_i\, \psi_i,
\]
which is equation~\eqref{eq:3line-psi}. Starting from any $\psi_i$ with non-zero overlap with the ground state of $\mathcal{H}_i$, the higher-eigenvalue components decay exponentially faster than the ground-state component, and after renormalisation $\psi_i$ relaxes to the lowest eigenfunction of $\mathcal{H}_i$ on $\Omega_i$, subject to the Bernoulli free-boundary condition $\partial\psi_i / \partial n = 0$ on $\Gamma_i$.

This is a standard technique: imaginary-time propagation is the variational method realised as a flow, with the energy strictly decreasing along the flow until the minimum is reached. The non-self-consistent piece of $\mathcal{H}_i$ depends only on the kernel potential $W$ and on the instantaneous Hartree potentials $V_j$ from the \emph{other} electrons; self-repulsion is excluded because $V_i$ itself does not appear in the sum.

\paragraph{\eqref{eq:3line-V} is the Poisson equation.} The Hartree potential $V_i$ generated by the unit electronic charge density $\psi_i^2$ satisfies, by Newton's identity for the inverse Laplacian in 3D,
\[
-\Delta V_i \;=\; 4\pi\, \psi_i^2 \quad\text{in } \mathbb{R}^3,
\]
with the convention $V_i(\mathbf{x}) = \int \psi_i^2(\mathbf{y}) / (2|\mathbf{x}-\mathbf{y}|) \,d\mathbf{y}$ from~\eqref{eq:PE} carrying a factor of $\tfrac{1}{2}$ that absorbs into the source term and gives the $2\pi$ of equation~\eqref{eq:3line-V}. Rather than inverting the Laplacian directly --- which would require a global linear solve at every step --- we relax to the Poisson solution by parabolic time-stepping:
\[
\dot{V}_i \;-\; \Delta V_i \;=\; 2\pi\, \psi_i^2,
\]
the linear heat equation with a source term. As $\dot{V}_i \to 0$, the equation reduces to $-\Delta V_i = -2\pi\psi_i^2$, the Poisson equation. The parabolic-relaxation form is what equation~\eqref{eq:3line-V} writes.

The advantage of relaxing rather than solving is that each step is a local stencil update: $V_i$ at each grid point depends only on its immediate neighbours. There is no global linear-algebra step, no FFT, and no boundary-condition machinery beyond the natural decay $V_i \to 0$ at infinity. The Poisson update is therefore one local kernel on the GPU per grid point per electron per time step.

\paragraph{\eqref{eq:3line-chi} is a level-set evolution.} The domain $\Omega_i$ is represented by an indicator function $\chi_i : \mathbb{R}^3 \to \{0, 1\}$ (or, in practice, a smoothed version of it). The boundary $\Gamma_i$ is the level set $\{\chi_i = \tfrac{1}{2}\}$. Equation~\eqref{eq:3line-chi} evolves the level set by advecting it with normal speed $\beta_i$:
\[
\dot{\chi}_i \;+\; \beta_i\, |\nabla \chi_i| \;=\; 0,
\]
the standard level-set transport equation~\cite{OsherSethian}. The speed $\beta_i$ at the inter-electron interface is set by the local jump of $\psi_i^2$ across the boundary, moving the front toward configurations satisfying the Bernoulli condition (continuity of $\psi_i$ at $\Gamma_i$). At convergence, $\dot{\chi}_i = 0$ and the partition is stationary.

Explicit time-stepping of \eqref{eq:3line-psi}--\eqref{eq:3line-chi} on a fixed real-space grid, with $\psi_i$ renormalised at every step, constitutes the entire solver. Each line is one PDE that becomes one stencil-update kernel on the GPU; everything else (forces on nuclei, dipole moments, Lindemann diagnostics, kernel softening, angular splitting) is post-processing or initialisation.

\section{Real-space grid and discretisation}\label{sec:grid}

The simulation domain is a cubic box of side $L$ atomic units, discretised as $N_{\rm g} \times N_{\rm g} \times N_{\rm g}$ grid points. Typical values are $N_{\rm g} = 100$--$200$, $L = 10$--$22$ a.u., grid spacing $h = L / N_{\rm g} \approx 0.1$ a.u. Each electron's density is represented as an $N_{\rm g}^3$ array of single-precision floats: the orbital amplitudes $\psi_i(\mathbf{x}_p)$ sampled at grid points $\mathbf{x}_p$. Per-electron Hartree potentials $V_i$ and indicator fields $\chi_i$ are stored on the same grid.

The Laplacian in~\eqref{eq:3line-psi} and~\eqref{eq:3line-V} is discretised by the standard 7-point stencil in 3D:
\[
\Delta f(\mathbf{x}_p) \;\approx\; \frac{1}{h^2}\Bigl(\sum_{q \sim p} f(\mathbf{x}_q) - 6 f(\mathbf{x}_p)\Bigr),
\]
where the sum is over the six nearest-neighbour grid points. The kinetic term, the Hartree-potential update, and the level-set transport are all local stencil operations of this kind. Each compute kernel touches a constant number of neighbours per grid point; the cost per time step is therefore $O(N_{\rm g}^3)$ in the number of grid points and $O(N_{\rm e})$ in the number of explicit electrons, giving a total cost of $O(N_{\rm e}\,N_{\rm g}^3)$ per step. This is the central reason RealQM is feasible on a laptop GPU at chemical accuracy: there are no global solves and no $O(N^k)$ matrix operations for $k>1$.

A typical step does approximately 10 GPU compute dispatches: orbital imaginary-time update, parabolic Poisson update, kinetic-energy reduction, normalisation, level-set update, force computation, and a few diagnostics. On a consumer laptop GPU with $\sim$10 TFLOPS, this runs at 30 or more steps per second for the systems reported in Part~III.

\paragraph{Box size and boundary effects.} The cubic box $L \times L \times L$ is large enough that the wave functions $\psi_i$ decay to numerical zero at the boundary. Typical bond-relevant configurations require $L \gtrsim 5$ a.u.\ around each atom in each direction; a 10-atom cluster with $\sim 4$ a.u.\ atomic spread needs $L \sim 16$--$20$ a.u. Boundary conditions are zero Dirichlet for both $\psi_i$ and $V_i$ on the outer faces of the box; for chemistry-relevant systems with bound electrons this is exact to numerical precision, since both quantities decay exponentially.

\section{Normalisation and self-interaction removal}\label{sec:norm-self}

Two ingredients of the iteration require explicit attention: the unit-norm constraint on each $\psi_i$ and the exclusion of self-repulsion in the Hartree sum.

\paragraph{Renormalisation.} Imaginary-time propagation~\eqref{eq:3line-psi} does not preserve the norm of $\psi_i$ automatically. After each step, $\psi_i$ is rescaled to satisfy $\int_{\Omega_i} \psi_i^2 = 1$:
\[
\psi_i(\mathbf{x}) \;\leftarrow\; \frac{\psi_i(\mathbf{x})}{\sqrt{h^3 \sum_p \psi_i(\mathbf{x}_p)^2}}.
\]
This is a single reduction over the grid points in $\Omega_i$ per electron per step. On the GPU the reduction is a standard parallel sum; the rescaling is then a pointwise multiplication.

\paragraph{Self-interaction.} The Hamiltonian $\mathcal{H}_i$ in~\eqref{eq:3line-psi} contains the sum $\sum_{j \ne i} V_j$ of Hartree potentials from the \emph{other} electrons, with $V_i$ itself excluded. In the implementation we store the total Hartree potential
\[
V_{\rm tot}(\mathbf{x}) \;=\; \sum_{j=1}^{N} V_j(\mathbf{x})
\]
on the grid, and form $V_{\rm tot} - V_i$ as the input to the $\mathcal{H}_i$ update for electron $i$. This is the standard self-interaction-correction operation, made trivial here because the per-electron potentials $V_i$ are explicitly maintained.

The self-interaction correction is exact in RealQM in a way it is not in approximate DFT: in DFT the exchange--correlation functional contains a residual self-repulsion that has to be subtracted heuristically, often imperfectly. In RealQM each $V_i$ is the exact Hartree potential of the unit density $\psi_i^2$, computed by Poisson relaxation~\eqref{eq:3line-V}, and the subtraction $V_{\rm tot} - V_i$ removes self-interaction exactly.

\section{The free-boundary update}\label{sec:bernoulli-update}

Equation~\eqref{eq:3line-chi} is the formal level-set evolution of the indicator $\chi_i$ of $\Omega_i$. In practice, two implementation variants of the same physical idea are used in the codebase.

\paragraph{Smoothed indicator.} In the simpler variant, $\chi_i$ is replaced by a smoothed indicator $w_i : \mathbb{R}^3 \to [0, 1]$ that softens the partition during the relaxation. The wave function on the smoothed partition is
\[
\psi(\mathbf{x}) \;=\; \sum_i w_i(\mathbf{x})\, \tilde\psi_i(\mathbf{x}),
\]
where $\tilde\psi_i$ is the unit-norm one-electron wave function and $w_i$ multiplies it by a partition-of-unity weight $\sum_i w_i(\mathbf{x}) = 1$. As the relaxation proceeds, $w_i$ tightens toward the indicator of $\Omega_i$: $w_i \to 1$ on $\Omega_i$ and $w_i \to 0$ off $\Omega_i$. At convergence, the smoothed partition matches the sharp Bernoulli partition to within the grid resolution.

The advantage of the smoothed variant is that derivatives across the interface are well-defined throughout the relaxation, and the kinetic-energy integral $\tfrac{1}{2} \int |\nabla \Psi|^2$ is computed without special handling at the boundary. The disadvantage is that the partition is not strictly non-overlapping during the iteration, which has to be reconciled at convergence.

\paragraph{Voronoi labelling.} In the second variant, $\Omega_i$ is represented by a labelling field $\ell(\mathbf{x}) \in \{1, \ldots, N\}$ that assigns each grid point to exactly one electron. The boundary $\Gamma_i$ is the discrete set of grid points where $\ell$ changes value. This is a Voronoi-like partition, with the boundary moved by reassigning grid points from one electron to another according to which $\psi_i^2$ is locally larger.

The labelling variant is strictly non-overlapping at every step, at the cost of having to handle the boundary specially when computing the kinetic energy near $\Gamma_i$. The two variants converge to the same Bernoulli partition in the continuum limit.

The choice between smoothed and labelling variants is made per implementation: \texttt{mol\_fast.js} uses unit-density orbitals with explicit angular splitting and a smoothed partition; \texttt{molecule.js} uses a Voronoi-partition labelling field. Both are described in Chapter~\ref{ch:webgpu}.

\section{Convergence diagnostics}\label{sec:convergence}

The relaxation is driven by three indicators, all monitored at every step:

\begin{itemize}
\item \emph{Total energy} $TE = K + PK + PE$. The variational principle guarantees that $TE$ decreases monotonically along the imaginary-time flow~\eqref{eq:3line-psi}, provided the Poisson and level-set updates have caught up with the wave-function update. In practice the three updates are interleaved at each time step, and $TE$ decreases monotonically except for small oscillations at the convergence floor.
\item \emph{Hartree residual} $\|\dot V_i\|$. The parabolic Poisson update~\eqref{eq:3line-V} has converged when $\dot V_i \to 0$. We monitor the maximum residual across electrons; convergence to numerical precision is reached when the residual falls below a fixed threshold.
\item \emph{Boundary motion} $\|\dot \chi_i\|$. The level-set update~\eqref{eq:3line-chi} has converged when the partition is stationary, $\dot \chi_i \to 0$. Boundary motion is the slowest of the three indicators to converge; in practice, the system is judged converged when boundary motion falls below grid resolution per step.
\end{itemize}

The three indicators are coupled: a partition update changes the domains over which the wave-function update operates, which in turn changes the Hartree sources for the Poisson update. The interleaved iteration converges typically in $10^2$--$10^4$ steps depending on system size and initial condition, corresponding to seconds to minutes of wall-clock time on a laptop GPU.

For molecular dynamics or geometry optimisation, the inner electronic relaxation is run to convergence at each fixed nuclear configuration; the nuclear positions are then updated by~\eqref{eq:force-coulomb}, and the inner relaxation resumes from the previous wave functions. The system is run as a coupled electron--nuclei dynamics in which the electrons are kept at their instantaneous variational minimum --- Born--Oppenheimer dynamics, realised through real-time relaxation rather than through energy gradients.

\section{Summary: what the algorithm is}\label{sec:algo-summary}

The full content of the RealQM solver, at the level of algorithm, is:

\begin{enumerate}
\item Initialise: place nuclei, initialise wave functions $\psi_i$ as suitable shell- or bond-shaped functions, initialise the partition $\{\Omega_i\}$.
\item Repeat until converged:
  \begin{enumerate}
  \item Update each $\psi_i$ by one imaginary-time step~\eqref{eq:3line-psi}.
  \item Renormalise each $\psi_i$ to unit norm on $\Omega_i$.
  \item Update each $V_i$ by one parabolic Poisson step~\eqref{eq:3line-V}.
  \item Update the partition by one level-set step~\eqref{eq:3line-chi}.
  \item Optionally compute forces on nuclei via~\eqref{eq:force-coulomb} and update nuclear positions.
  \end{enumerate}
\item Report total energy, geometry, dipole moments, and any other diagnostics.
\end{enumerate}

That is the algorithm. The variational character of the problem, the non-overlap geometry, the free-boundary condition, and the local-potential force law are all built into these steps. Chapter~\ref{ch:webgpu} describes how these steps are mapped onto a GPU compute pipeline that runs interactively in a browser.

% --- End of Chapter 4 ---
